% ---------------------------------------------------------------
% Research audit of "Visual DOM-Graph RAG: Online Evidence Control
% for Multimodal Web-Page RAG" (A. Firome)
% Target venue: DocInsights 2026 @ EMNLP 2026, Budapest
% Compile: pdflatex domgraph_audit.tex  (run twice for refs)
% ---------------------------------------------------------------
\documentclass[11pt,a4paper]{article}

\usepackage[T1]{fontenc}
\usepackage[utf8]{inputenc}
\usepackage[margin=2.4cm]{geometry}
\usepackage{booktabs}
\usepackage{tabularx}
\usepackage{longtable}
\usepackage{array}
\usepackage{enumitem}
\usepackage{xcolor}
\usepackage{amssymb}
\usepackage{titlesec}
\usepackage[colorlinks=true,linkcolor=black,urlcolor=blue!60!black,citecolor=black]{hyperref}

% ---- severity tags -------------------------------------------------
\definecolor{blockred}{HTML}{A61B29}
\definecolor{majororange}{HTML}{B25E00}
\definecolor{modblue}{HTML}{1F5C8B}
\definecolor{mingray}{HTML}{555555}
\definecolor{okgreen}{HTML}{1E6B3A}

\newcommand{\blocking}{\textcolor{blockred}{\textbf{BLOCKING}}}
\newcommand{\major}{\textcolor{majororange}{\textbf{MAJOR}}}
\newcommand{\moderate}{\textcolor{modblue}{\textbf{MODERATE}}}
\newcommand{\minor}{\textcolor{mingray}{\textbf{MINOR}}}
\newcommand{\ok}{\textcolor{okgreen}{\textbf{OK}}}
\newcommand{\finding}[2]{\subsection*{#1 \hfill #2}\addcontentsline{toc}{subsection}{#1}}

\newcolumntype{L}[1]{>{\raggedright\arraybackslash}p{#1}}

\titleformat{\section}{\large\bfseries}{\thesection}{0.7em}{}
\setlist[itemize]{itemsep=2pt,topsep=3pt,leftmargin=1.4em}
\setlist[enumerate]{itemsep=2pt,topsep=3pt,leftmargin=1.6em}

\title{\textbf{Research Audit}\\[0.35em]
\large ``Visual DOM-Graph RAG: Online Evidence Control for\\ Multimodal Web-Page RAG''}
\author{Prepared for submission planning --- DocInsights 2026 (EMNLP 2026, Budapest)}
\date{30 July 2026}

\begin{document}
\maketitle

\begin{center}
\begin{tabularx}{\textwidth}{@{}lX@{}}
\toprule
\textbf{Manuscript} & Visual DOM-Graph RAG: Online Evidence Control for Multimodal Web-Page RAG \\
\textbf{Author} & Alexis Firome (single author, no institutional affiliation listed) \\
\textbf{Length} & 8 pages incl. references; no Limitations or Ethics section \\
\textbf{Target} & DocInsights 2026 --- Workshop on Document Intelligence and Understanding \\
\textbf{Direct deadline} & 2 August 2026, 23:59 AoE \hfill \textcolor{blockred}{(3 days from this audit)} \\
\textbf{ARR commitment} & 30 August 2026 (requires completed ARR reviews + meta-review) \\
\bottomrule
\end{tabularx}
\end{center}

\vspace{0.6em}

% =====================================================================
\section{Headline verdict}

The manuscript describes a coherent, well-motivated, and apparently fully implemented
system. Its engineering narrative is unusually honest --- the oversized-page splitting
design is presented together with the failure that motivated it, and the author explicitly
refuses to report numbers that have not been produced. That honesty is a genuine strength
and should be preserved.

However, \textbf{the paper is not currently submittable as an archival full paper}. It
contains no experimental results of any kind, at least eight unresolved cross-references
(\verb|\S??|), placeholder text inside the bibliography, figures captioned in French, no
Limitations section, and an evaluation protocol that, as written, cannot be executed on the
benchmarks it names. Several of these are desk-reject triggers at ACL-affiliated venues
independent of scientific merit.

\begin{center}
\fcolorbox{blockred}{blockred!5}{%
\begin{minipage}{0.95\textwidth}
\vspace{2pt}
\textbf{Recommendation.} Submit to the \textbf{direct, non-archival} track by 2 August 2026
as a work-in-progress system description aimed at the poster/demo session, after completing
the four-item triage list in \S\ref{sec:plan}. Do \emph{not} submit as archival: an archival
paper with an empty results section will be scored against a full-paper rubric and is very
likely to be rejected, which would also burn the venue. If a genuine evaluation can be
completed within four weeks, the stronger play is a full archival submission to a later
venue rather than an ARR commitment here (the 30 August ARR route requires reviews that are
already complete, which is not the case).
\vspace{3pt}
\end{minipage}}
\end{center}

% =====================================================================
\section{Fit with the DocInsights 2026 call}

Topical fit is strong; the mismatch is with the venue's expectations of evidence, not its scope.

\begin{center}
\begin{tabularx}{\textwidth}{@{}L{5.2cm}cX@{}}
\toprule
\textbf{Workshop theme} & \textbf{Fit} & \textbf{Comment} \\
\midrule
Structure-aware modeling & \ok & Direct hit. Layout, reading order, hierarchy and cross-element relations are the paper's core subject. \\
Multimodal document understanding & \ok & Joint text--layout--vision; OCR-free reading of rendered tiles. \\
Long and cross-document intelligence & \ok & The \texttt{continues} and \texttt{links\_to} relations address exactly the multi-page/collection framing. \\
Reasoning and grounding & \moderate & Evidence selection is central, but attribution, verification and hallucination mitigation are not measured. \\
Evaluation and robustness & \blocking & The workshop explicitly solicits benchmarks, diagnostics and failure-mode analysis. This paper contributes none, and this is the axis reviewers will weigh hardest. \\
Deployment and governance & \minor & Not addressed; not required. \\
\bottomrule
\end{tabularx}
\end{center}

\noindent The CFP language (``systems, demos, diagnostic studies, and work in progress''
in the poster session) is unusually accommodating of exactly this manuscript's maturity ---
but only through the non-archival door.

% =====================================================================
\section{Blocking issues}

\finding{B1. The evaluation section does not exist}{\blocking}
The abstract, the contributions list, the related-work section and four separate method
subsections all forward-reference results to ``\S??''. There is no \S5, no results table,
no qualitative error analysis, and no baseline comparison. Every empirical claim in the
paper is a promise. Reviewers cannot assess a system whose only reported behaviour is that
it runs.

\emph{Minimum viable fix within 3 days:} a single-page retrieval evaluation on a held-out
slice of the existing synthetic QA set --- Recall@$k$ and MRR for (a) DOM-guided tiles vs.\
(b) a fixed-height pixel grid reimplemented over the same screenshots, with the same
untuned scorer. This is a few hours of work with the code already written, it directly
tests the paper's central claim, and it converts the submission from ``no results'' to
``preliminary results''.

\finding{B2. The named benchmarks cannot run this pipeline}{\blocking}
The stated protocol is comparison against MAGE-RAG's published LongDocURL and
MMLongBench-Doc numbers. Both are \emph{PDF} long-document benchmarks. The entire method
is predicated on the availability of a live DOM obtained by rendering an HTML page in
headless Chromium: bounding boxes come from \texttt{getBoundingClientRect}, hyperlinks come
from \texttt{<p>} elements, chrome removal uses CSS selectors, and page splitting uses
\texttt{h1}/\texttt{h2} boundaries. None of these inputs exist for a scanned or
born-digital PDF. As written, the evaluation protocol is not executable.

This is the most damaging technical criticism available to a reviewer, and it is visible
from the abstract alone. It requires an explicit answer, not a hedge. Options, in
decreasing order of strength:
\begin{enumerate}
\item Evaluate on an HTML-native benchmark instead. Candidates worth checking:
      WebSRC, WebQA, MMLongBench-Doc's web-sourced subsets if any, or a purpose-built
      Wikipedia QA set with human-verified answers.
\item Construct and release a small HTML document-QA benchmark as a contribution in its
      own right --- this would fit the DocInsights ``benchmarks'' theme and turn a
      weakness into a second contribution.
\item Define a PDF$\rightarrow$HTML ingestion path and evaluate the degradation, conceding
      that the DOM is then reconstructed rather than authoritative --- which weakens the
      paper's core argument and should be a last resort.
\end{enumerate}
Additionally, comparing against \emph{published} numbers rather than re-running baselines
is not a controlled comparison: reader model, action budget, prompt and retrieval depth all
differ. State this explicitly or re-run.

\finding{B3. Venue formatting and anonymity non-compliance}{\blocking}
\begin{itemize}
\item The manuscript is not in the ACL style. ACL-affiliated workshops require the official
      ACL template (typically 8 pages long / 4 pages short, references unlimited). Verify
      the exact limits against the DocInsights CFP before reformatting.
\item \textbf{No Limitations section.} This is mandatory under ACL policy for all
      submissions, does not count against the page limit, and its absence is a stated
      desk-reject condition at ACL-family venues. Given this paper's status, the Limitations
      section is also the natural and dignified home for the missing-results discussion.
\item \textbf{Anonymity.} The title page carries the author's name and a personal Gmail
      address. If review is anonymous, this alone is a desk reject. The accompanying
      notebook, batch script and ``open-sourced'' repository must be anonymised (e.g.\ an
      anonymous.4open.science mirror) for the review period.
\item No Ethics statement and no Responsible NLP / reproducibility checklist.
\end{itemize}

\finding{B4. Unresolved cross-references and bibliography placeholders}{\blocking}
Eight or more \verb|\S??| markers survive into the compiled PDF, including one in the
abstract. Five bibliography entries carry editorial notes to the author --- ``Author list
to verify precisely before final submission'' appears in refs [3], [6], [7], [8] and [9].
These must not reach a reviewer. Verify every author list and venue against the primary
sources; two of the five (LongDocURL, MMLongBench-Doc) are published and easy to fix.

\finding{B5. Figures are captioned in French in an English paper}{\blocking}
Figures 2, 3 and 4 carry in-image titles and legends in French (``Tuiles et ordre de
lecture, superpos\'es \`a la page''; ``Graphe page/\'el\'ements --- x = position horizontale
r\'eelle''; ``Vue page-de-pages''; category legends ``Texte / Tableau / Liste''). Regenerate
the plots with English strings. Trivial to fix, immediately noticed, and it signals
carelessness to a reviewer before they read a word of the method.

% =====================================================================
\section{Major technical concerns}

\finding{M1. The synthetic training data is constructed so that the contribution cannot help}{\major}
QA generation prompts the VLM for a question ``whose answer is visible \emph{only} in that
tile.'' Every training and (presumably) evaluation example is therefore single-hop and
single-tile by construction. But the graph and the evidence controller exist precisely to
support evidence that is distributed across tiles, columns, sections and pages. On this
data distribution, an oracle system opens exactly one node, the frontier expansion along
\texttt{reading\_order} can never be necessary, and \texttt{layout\_adjacency},
\texttt{section\_hierarchy}, \texttt{continues} and \texttt{links\_to} are provably inert.
The experiment, as designed, cannot demonstrate the paper's thesis. A multi-tile /
multi-hop question generation mode is needed --- e.g.\ prompt over a pair of adjacent or
hyperlinked tiles and require both to be opened.

\finding{M2. The controller is a lexical text retriever wearing a vision system's clothes}{\major}
Node relevance is scored by TF-IDF cosine similarity between the query and the node's
\emph{extracted DOM text}. Hard negatives are mined the same way. The paper's opening
argument is that text extraction discards the information that layout encodes --- yet at
query time, which tiles the VLM is ever allowed to see is decided entirely by a bag-of-words
match over extracted text. The visual reader can only rerank what a lexical matcher has
already shortlisted, so any reported end-to-end number would be upper-bounded by lexical
recall and would not test the visual pathway at all.

The paper acknowledges this as a ``deliberate placeholder,'' which is honest but not
sufficient: the placeholder sits on the critical path of every claim. Wiring the fine-tuned
reader's question--tile similarity into the controller is the single highest-value remaining
piece of engineering, and until it is done the system's headline positioning
(``OCR-free, visual'') is not accurate.

\finding{M3. The central claim --- DOM tiling beats grid tiling --- is never tested}{\major}
The motivating failure mode (a table sliced mid-row) is asserted, illustrated once
qualitatively on a short Wikipedia biography, and never quantified. Nothing in the paper
establishes that content-aligned tiles improve retrieval, and there is a plausible
counter-hypothesis a reviewer will raise: fixed-height grids produce overlapping context
and uniform token cost, whereas DOM tiles vary wildly in aspect ratio and may be worse for
a VLM encoder with a fixed patch budget. Report tiles/page, pixels/tile, and token cost
alongside accuracy for both tilings.

\finding{M4. Self-distillation and evaluator--generator coupling}{\major}
Questions are generated by a VLM (\texttt{qwen3-vl} via Ollama, per \S4), the optional
false-negative judge is the same reader VLM, and the fine-tuned model is
\texttt{Qwen2-VL-7B-Instruct}. Note in passing that the generator and the fine-tuned reader
are different model generations, which is not flagged in the text and should be. More
importantly: training data, negative filtering and (if used for eval) ground truth all
originate from the same model family. Any gain may reflect distillation of the generator's
idiosyncrasies rather than improved document understanding. At minimum, hand-verify a
random sample of $\sim$100 generated pairs and report the error rate; ideally evaluate on
human-authored questions.

\finding{M5. Four of six graph relations are never consumed}{\major}
Algorithm 1 expands the frontier along \texttt{reading\_order} neighbours only. The paper
concedes the controller does not cross \texttt{continues} edges. \texttt{layout\_adjacency}
and \texttt{section\_hierarchy} are constructed, visualised, and then unused. The
multigranular graph is therefore, operationally, a linear chain plus decoration. Either
extend the controller to type-aware expansion (with per-relation ablations --- an excellent,
cheap experiment that would substantially strengthen the paper) or drop the unused relations
and reduce the claim.

\finding{M6. Wikipedia-only, with Wikipedia-specific hardcoding}{\major}
The chrome-removal selector list (\texttt{.navbox}, \texttt{.mw-editsection}, \dots) is
MediaWiki-specific, and Wikipedia's HTML is unusually clean and semantically marked up.
The generalisation claim is about ``rendered web pages,'' but news sites, SPAs,
JavaScript-hydrated content, cookie walls and \texttt{div}-soup layouts will degrade DOM
quality badly. Either evaluate on at least one non-Wikipedia domain or explicitly scope the
claim to semantically marked-up encyclopaedic content and move the general case to
Limitations.

\finding{M7. SimpleQA is the wrong instrument}{\moderate}
SimpleQA measures parametric short-form factuality with no supplied documents. It cannot
serve as a sanity check for a document-grounded retrieval pipeline --- there is no page to
tile, no graph to build, and no evidence to control. Remove it, or replace it with a
closed-book/open-book contrast that actually isolates whether evidence-controlled reading
helps or hurts on questions the model already knows.

\finding{M8. Every hyperparameter is unmotivated}{\moderate}
80\,px minimum tile height, 2048\,px maximum tile height and viewport width, 20\,px
vertical-overlap tolerance, $k{=}3$ link-bearing tiles, $n{=}2$ negatives, a $10\times$
oversize threshold, 0.5 scale factor, $\tau{=}0.05$, LoRA $r{=}16$, $\alpha{=}32$. No
sensitivity analysis, no justification, no indication of how any was chosen. Absolute pixel
thresholds tied to a single 2048\,px viewport are also a portability hazard for responsive
layouts and high-DPI rendering. At least one sensitivity curve (minimum tile height vs.\
retrieval quality) would be persuasive and is cheap.

\finding{M9. Dataset and cost are unreported}{\moderate}
The paper never states how many pages, sections, tiles, QA pairs or mined negatives exist;
what fraction of tiles returned SKIP; how often the lexical filter and the VLM judge fired;
training time; GPU type; or inference cost. A reader cannot judge whether this is 30 tiles
or 30{,}000. A dataset-statistics table is mandatory and takes minutes to produce from
\texttt{contrastive\_examples.jsonl}.

\finding{M10. An unattributed direct quotation}{\major}
\S1 quotes a criticism of PixelRAG's tiling verbatim (``slices pages by fixed pixel
height\dots'') and attributes it only to ``independent coverage of the method,'' with no
citation. A verbatim quotation with no source is a research-integrity problem regardless of
how minor the source is. Either cite the source properly, or --- better --- state the
limitation in your own words with a pointer to the relevant section of the PixelRAG paper
or its repository, where the fixed grid is presumably visible in the code.

% =====================================================================
\section{Claim--evidence ledger}

Every substantive claim in the manuscript, and what currently backs it.

\begin{center}
\small
\begin{longtable}{@{}L{7.1cm}L{4.1cm}L{4.4cm}@{}}
\toprule
\textbf{Claim} & \textbf{Evidence provided} & \textbf{Status} \\
\midrule
\endhead
DOM-guided tiles follow content boundaries by construction & Method description; Fig.~2 on one page & \textcolor{okgreen}{Supported} (definitional) \\
\dots and therefore improve retrieval over a pixel grid & None & \textcolor{blockred}{Unsupported} \\
Nested-element pruning avoids double-tiling & Unit tests & \textcolor{okgreen}{Supported} \\
Heading pairing improves answerability & None & \textcolor{blockred}{Unsupported} \\
Oversized-page splitting preserves long-document coverage & Anecdote (46{,}000\,px page); Fig.~4 & \textcolor{majororange}{Partial} --- feasibility only \\
Contrastive LoRA yields a better tile reader & None; training implemented, not run to completion & \textcolor{blockred}{Unsupported} \\
Hard negatives are useful contrastive signal & Argued a priori & \textcolor{blockred}{Unsupported} \\
VLM-judge filter catches non-lexical false negatives & Stub-client test only & \textcolor{blockred}{Unsupported} \\
Reduced resolution preserves accuracy at lower cost & Infrastructure only; runs pending & \textcolor{blockred}{Unsupported} \\
Budgeted best-first control beats static top-$k$ & Argued; ablation deferred & \textcolor{blockred}{Unsupported} \\
Graph relations beyond reading order add value & None; unused by the controller & \textcolor{blockred}{Unsupported} \\
Pipeline runs end-to-end on live pages & Notebook, batch script, smoke test & \textcolor{okgreen}{Supported} \\
\bottomrule
\end{longtable}
\end{center}

\noindent Three of twelve claims are supported, and the three that are supported are the
implementation claims rather than the scientific ones. This ratio is the audit's core
finding.

% =====================================================================
\section{Related-work gaps}

The related-work section rests on two 2026 arXiv preprints (PixelRAG, MAGE-RAG), neither
peer-reviewed, plus standard dense-retrieval references. For a DocInsights audience, the
visual document retrieval literature is conspicuously absent. A reviewer working in this
area will notice immediately. Verify current status and add, at minimum:

\begin{itemize}
\item \textbf{Screenshot / late-interaction visual retrieval:} ColPali and ColQwen-style
      page-image retrievers, and Document Screenshot Embedding (DSE). These are the
      canonical prior art for ``retrieve over rendered pages with a VLM,'' and their absence
      makes PixelRAG look more novel than it is --- which weakens \emph{your} positioning
      too, since your contribution is orthogonal to theirs and could be composed with them.
\item \textbf{Visual RAG pipelines:} VisRAG, M3DocRAG, and related multi-page visual RAG systems.
\item \textbf{Graph-structured RAG:} GraphRAG and its successors, to situate the
      graph-plus-controller design in a broader lineage than a single preprint.
\item \textbf{Layout- and markup-aware pretraining:} the LayoutLM family, MarkupLM, and
      DOM/HTML-structured models --- directly relevant, since you argue the DOM is a
      privileged signal.
\item \textbf{DOM + screenshot agents:} WebArena, Mind2Web, SeeAct, and accessibility-tree
      approaches. These communities have already studied when DOM structure is reliable and
      when it is not, which bears on M6.
\item \textbf{VLMs as universal encoders:} the ``last-token hidden state as embedding''
      technique is used in \S3.3 and attributed only to ``recent work.'' Cite it.
\end{itemize}

\noindent Two further uncited assertions to fix: ``a broader trend of using VLMs as
universal, modality-agnostic retrievers'' (\S2) and ``Pillow's recommended high-quality
downscale filter'' (\S3.3).

% =====================================================================
\section{Compliance checklist}

\begin{center}
\begin{tabularx}{\textwidth}{@{}L{6.6cm}cX@{}}
\toprule
\textbf{Requirement} & \textbf{Status} & \textbf{Action} \\
\midrule
Official ACL style template & \textcolor{blockred}{No} & Port to the ACL template; verify version against the CFP. \\
Page limit (verify: typically 8 long / 4 short) & \textcolor{majororange}{Unknown} & Currently 8pp \emph{including} references, which suggests headroom once references are excluded. \\
Limitations section (mandatory, uncounted) & \textcolor{blockred}{Missing} & Write it. Cover: no results, Wikipedia-only, lexical scorer, single-hop data, PDF benchmarks inapplicable. \\
Ethics statement & \textcolor{majororange}{Missing} & Short statement: crawling, licensing, no human subjects. \\
Anonymisation (if double-blind) & \textcolor{blockred}{No} & Remove name/email; anonymise repo; check the PDF metadata too. \\
Responsible NLP / reproducibility checklist & \textcolor{majororange}{Missing} & Complete at submission time. \\
Resolved cross-references & \textcolor{blockred}{No} & 8+ instances of \verb|\S??|. \\
Clean bibliography & \textcolor{blockred}{No} & Remove five ``verify before submission'' notes. \\
Consistent system name & \textcolor{mingray}{No} & Title says ``Visual DOM-Graph RAG''; body says ``DOM-Graph RAG''. Pick one. \\
Figures legible and in English & \textcolor{blockred}{No} & Regenerate Figs.~2--4; Fig.~4's rotated overlapping labels are also unreadable at print size. \\
Archival vs.\ non-archival declared & --- & Choose non-archival (see \S1). \\
\bottomrule
\end{tabularx}
\end{center}

% =====================================================================
\section{Reproducibility and integrity}

\textbf{Strengths.} The code-level reproducibility story is better than most workshop
submissions: one module per pipeline stage, dependency-free unit tests over synthetic
fixtures, a Playwright-only end-to-end smoke test that needs no GPU or paid API, LoRA target
modules discovered by introspection rather than hardcoded, per-section checkpointing to
\texttt{contrastive\_examples.jsonl}, and a resolution ablation designed so that only one
variable changes between runs. The manifest-stores-paths-not-pixels design is a genuinely
neat detail worth keeping prominent.

\textbf{Gaps.}
\begin{itemize}
\item No random seeds, no repeated runs, no variance reporting --- and with a small
      synthetic set, single-run differences will not be meaningful. Commit to $\geq 3$ seeds
      before any number is reported.
\item Live Wikipedia is a moving target. Pages change; the reported dataset will not be
      reconstructible in six months. Archive the rendered screenshots and the extracted DOM
      geometry (or pin \texttt{oldid} revision IDs) and release them, not just the crawler.
\item Rendering is browser-version-dependent: \texttt{getBoundingClientRect} output varies
      with Chromium version, fonts installed, and device pixel ratio. Pin and report the
      Playwright/Chromium version and the font environment (ideally ship a container).
\item Model versions are underspecified: which \texttt{qwen3-vl} Ollama tag, which
      \texttt{Qwen2-VL-7B-Instruct} revision.
\item No licence stated for the released code or the derived tile images.
\end{itemize}

% =====================================================================
\section{Ethics, licensing and data governance}

Low risk overall, but three items belong in a short Ethics statement:
\begin{itemize}
\item \textbf{Crawling.} The corpus builder fetches up to $m$ linked pages with retry and
      backoff. State that crawling respects \texttt{robots.txt} and rate limits, and prefer
      the Wikimedia REST API or a dump over live scraping where possible.
\item \textbf{Licensing.} Wikipedia text is CC BY-SA; embedded images carry heterogeneous
      licences, some non-free. Rendered screenshots and cropped tiles are derivative works.
      If the tile dataset is released, the licensing must be addressed explicitly --- this is
      the most likely real-world obstacle to releasing the dataset at all, and it is worth
      resolving before you promise a release in print.
\item \textbf{Generated content.} Synthetic questions from a VLM may contain errors or
      encode the model's biases; the release should label them as machine-generated and
      unverified.
\end{itemize}

% =====================================================================
\section{Prioritized action plan}
\label{sec:plan}

\subsection*{Before 2 August (3 days) --- must do}
\begin{enumerate}
\item Resolve every \verb|\S??|; strip the five bibliography placeholder notes; regenerate
      Figures 2--4 with English labels. \emph{(2--3 hours, pure hygiene, highest
      value-per-minute in this document.)}
\item Port to the ACL template; add a Limitations section and a short Ethics statement;
      anonymise author block, repository links and PDF metadata. \emph{(3--4 hours.)}
\item Rewrite the abstract and \S1 so the contribution is stated as \emph{a pipeline and a
      released implementation}, with evaluation declared as in progress in the Limitations
      section rather than promised via a broken cross-reference. Remove SimpleQA. Replace the
      LongDocURL/MMLongBench-Doc protocol with an honest paragraph acknowledging the
      PDF-vs-DOM mismatch (B2) and naming an HTML-native alternative. \emph{(3 hours ---
      this is the paragraph that decides the review.)}
\item If at all possible, run the minimal DOM-tiling vs.\ grid-tiling retrieval comparison
      from B1 and report it, however small, with the dataset-statistics table from M9.
      One real table changes the paper's category. \emph{(half a day.)}
\end{enumerate}

\subsection*{Before camera-ready (by 27 September, if accepted)}
\begin{enumerate}
\item Complete the contrastive fine-tuning runs at both resolutions; report the token-cost
      /accuracy tradeoff as an actual measurement rather than a restatement of PixelRAG's claim.
\item Replace the TF-IDF controller scorer with the fine-tuned reader's similarity (M2), and
      report both --- the lexical version becomes a legitimate ablation instead of a placeholder.
\item Add multi-tile question generation (M1) and a controller ablation: static top-$k$ vs.\
      best-first, and \texttt{reading\_order}-only vs.\ type-aware expansion (M5).
\item Add at least one non-Wikipedia domain (M6).
\end{enumerate}

\subsection*{For a future full paper}
Human-verified evaluation set; sensitivity analysis on the tiling thresholds; comparison
against a ColPali/DSE-class visual retriever as the real baseline rather than a
reimplemented grid; multi-seed results with variance.

% =====================================================================
\section{Anticipated reviewer questions}

Prepare answers to these; each is a likely rebuttal item.
\begin{enumerate}
\item How can this be compared to MAGE-RAG on PDF benchmarks when the method requires a DOM?
\item If node selection is TF-IDF over extracted text, in what sense is the system OCR-free?
\item Why should tiles that follow element boundaries help a VLM encoder that resamples to a
      fixed patch grid anyway?
\item Since every training question is answerable from a single tile, what work is the graph doing?
\item How does this differ from ColPali/DSE plus a chunker?
\item What is the tiling overhead per page relative to a fixed grid, and does it pay for itself?
\item What happens on a page whose DOM does not reflect the visual layout (CSS grid, absolute
      positioning, JS-injected content)?
\end{enumerate}

% =====================================================================
\appendix
\section{Inventory of unresolved references}

\begin{center}
\begin{tabularx}{\textwidth}{@{}lX@{}}
\toprule
\textbf{Location} & \textbf{Broken reference} \\
\midrule
Abstract & ``reported as a work-in-progress in Section ??'' --- an abstract should not carry a cross-reference at all \\
\S1, contributions preamble & ``\S?? lays out the intended evaluation'' \\
\S2, graph paragraph & ``only two of MAGE-RAG's five relation types so far (\S??)'' --- note this also contradicts \S3.2, which claims \emph{four} of five \\
\S2, factuality & ``we include it in our evaluation protocol (\S??)'' \\
\S3.2, missing relation & ``requires an embedding model; see \S??'' \\
\S3.2, entry points & ``does not yet cross \texttt{continues} edges automatically (\S??)'' \\
\S3.3, resolution & ``results are reported in \S?? once both runs complete'' \\
\S3.4, scoring & ``score nodes with its learned question--tile similarity instead (\S??)'' \\
\S3.4, policy & ``an ablation we report in \S??'' \\
\bottomrule
\end{tabularx}
\end{center}

\noindent The two-vs-four relation-count contradiction between \S2 and \S3.2 is a
substantive inconsistency, not just a stale reference --- fix the text, not only the label.

\section{Undocumented constants requiring justification}

\begin{center}
\begin{tabularx}{\textwidth}{@{}L{5.2cm}L{2.6cm}X@{}}
\toprule
\textbf{Constant} & \textbf{Value} & \textbf{Concern} \\
\midrule
Viewport width & 2048\,px & Ties all downstream thresholds to one rendering condition \\
Max tile height & 2048\,px & Square cap asserted, not derived from the encoder's patch budget \\
Min tile height (merge) & 80\,px & Unjustified; likely the most sensitive tiling parameter \\
Row-grouping tolerance & 20\,px & Will misgroup dense or high-DPI layouts \\
Nesting tolerance & 1\,px & Reasonable; state that it absorbs subpixel rounding \\
Link-bearing tiles kept & $k{=}3$ & Arbitrary crawl-branching bound \\
Mined negatives & $n{=}2$ & No sweep over the number of negatives \\
Oversize page threshold & $10\times$ max tile & Chosen post hoc from one pathological page \\
Downscale factor & 0.5 & Only one point on the resolution curve \\
InfoNCE temperature & $\tau{=}0.05$ & Conventional but unswept \\
LoRA rank / alpha / dropout & 16 / 32 / 0.05 & Conventional but unswept \\
\bottomrule
\end{tabularx}
\end{center}

\vfill
\begin{center}\footnotesize
\emph{This audit reflects the manuscript as supplied and the DocInsights 2026 CFP as
published. Page limits, anonymity policy and template version should be confirmed against
the workshop's OpenReview instructions, which are the system of record.}
\end{center}

\end{document}